% Part: first-order-logic
% Chapter: sequent-calculus
% Section: derivations

\documentclass[../../../include/open-logic-section]{subfiles}

\begin{document}

\iftag{FOL}
      {\olfileid{fol}{seq}{der}}
      {\olfileid{pl}{seq}{der}}

\olsection{\usetoken{P}{derivation}}

\begin{explain}
گفتیم دنبالهٔ آغازین چه صورتی دارد و قواعد استنتاج را نیز به دست دادیم.
!!^{derivation}s حساب دنباله‌ای به‌طور استقرایی از این‌ها ساخته می‌شوند:
هر !!{derivation} یا به‌تنهایی دنباله‌ای آغازین است، یا از
!!{derivation}s تشکیل می‌شود---یک یا دو مورد---که استنتاجی در پی دارند.
\end{explain}

\begin{defn}[$\Log{LK}$ !!{derivation}]
\emph{!!{derivation} در $\Log{LK}$} برای دنبالهٔ~$S$، درختی متناهی
از دنباله‌هاست که شرایط زیر را برآورده می‌کند:
\begin{enumerate}
\item بالاترین دنباله‌های درخت دنباله‌های آغازین‌اند.
\item پایین‌ترین دنبالهٔ درخت~$S$ است.
\item هر دنبالهٔ درخت جز $S$، مقدمهٔ کاربرد درستی از یک قاعدهٔ استنتاج
  است که نتیجهٔ آن بی‌درنگ زیر آن دنباله در درخت قرار دارد.
\end{enumerate}
سپس می‌گوییم $S$ \emph{دنبالهٔ پایانیِ} !!{derivation} است و
$S$ \emph{در $\Log{LK}$ !!{derivable} است} (یا در
$\Log{LK}$-!!{derivable} است).
\end{defn}

\begin{ex}
هر دنبالهٔ آغازین، برای نمونه $!C \Sequent !C$، !!a{derivation} است.
می‌توانیم با اعمال مثلاً قاعدهٔ $\LeftR{\Weakening}$،
\begin{prooftree}
\Axiom$ \Gamma \fCenter \Delta $
\RightLabel{\LeftR{\Weakening}}
\UnaryInf$ !A, \Gamma \fCenter \Delta$
\end{prooftree}
!!{derivation} تازه‌ای از آن به دست آوریم. بااین‌حال، مراد این است که
قاعده عام باشد: می‌توان $!A$ را در قاعده با هر !!{sentence}، برای نمونه
با~$!D$، جایگزین کرد. اگر مقدمه با دنبالهٔ آغازینِ
$!C \Sequent !C$ منطبق باشد، این بدان معناست که هم $\Gamma$ و هم
$\Delta$ صرفاً~$!C$ هستند و آنگاه نتیجه $!D, !C \Sequent !C$ خواهد
بود. بنابراین عبارت زیر !!a{derivation} است:
\begin{prooftree}
\Axiom$ !C \fCenter !C $
\RightLabel{\LeftR{\Weakening}}
\UnaryInf$ !D, !C \fCenter !C$
\end{prooftree}
اکنون می‌توانیم قاعده‌ای دیگر، مثلاً $\LeftR{\Exchange}$، را اعمال
کنیم که اجازه می‌دهد در میان !!{sentence}s سمت چپ، دو مورد را جابه‌جا
کنیم. پس عبارت
زیر نیز !!{derivation} درستی است:
\begin{prooftree}
\Axiom$ !C \fCenter !C $
\RightLabel{\LeftR{\Weakening}}
\UnaryInf$ !D, !C \fCenter !C$
\RightLabel{\LeftR{\Exchange}}
\UnaryInf$ !C, !D \fCenter !C$
\end{prooftree}
در این کاربرد قاعده که به صورت زیر داده شده بود،
\begin{prooftree}
\Axiom$ \Gamma, !A, !B, \Pi \fCenter \Delta $
\RightLabel{\LeftR{\Exchange}}
\UnaryInf$ \Gamma, !B, !A, \Pi \fCenter \Delta,$
\end{prooftree}
هم $\Gamma$ و هم $\Pi$ تهی بودند، $\Delta$ همان $!C$ است و نقش‌های
$!A$ و $!B$ را به‌ترتیب $!D$ و~$!C$ بازی می‌کنند. تقریباً به همین
شیوه می‌بینیم که
\begin{prooftree}
\Axiom$ !D \fCenter !D $
\RightLabel{\LeftR{\Weakening}}
\UnaryInf$ !C, !D \fCenter !D$
\end{prooftree}
نیز !!a{derivation} است. اکنون می‌توانیم این !!{derivation}s، یعنی
هر دو را،
برداریم و با استفاده از $\RightR{\land}$ ترکیب کنیم. آن قاعده چنین بود:
\begin{prooftree}
\Axiom$\Gamma \fCenter \Delta, !A$
\Axiom$ \Gamma \fCenter \Delta, !B$
\RightLabel{\RightR{\land}}
\BinaryInf$ \Gamma \fCenter \Delta, !A \land !B $
\end{prooftree}
در مورد ما، مقدمه‌ها باید با آخرین دنباله‌های !!{derivation}sی که به
آن مقدمه‌ها ختم می‌شوند منطبق باشند. یعنی $\Gamma$ همان $!C, !D$،
$\Delta$ تهی، $!A$ همان $!C$ و $!B$ همان $!D$ است. پس اگر استنتاج
درست باشد، نتیجه $!C, !D \Sequent !C \land !D$ است.
\begin{prooftree}
\Axiom$ !C \fCenter !C $
\RightLabel{\LeftR{\Weakening}}
\UnaryInf$ !D, !C \fCenter !C$
\RightLabel{\LeftR{\Exchange}}
\UnaryInf$ !C, !D \fCenter !C$
\Axiom$ !D \fCenter !D $
\RightLabel{\LeftR{\Weakening}}
\UnaryInf$ !C, !D \fCenter !D$
\RightLabel{\RightR{\land}}
\BinaryInf$ !C, !D \fCenter !C \land !D $
\end{prooftree}
البته می‌توانیم ترتیب مقدمه‌ها را نیز وارونه کنیم؛ آنگاه $!A$ همان
$!D$ و $!B$ همان~$!C$ خواهد بود.
\begin{prooftree}
\Axiom$ !D \fCenter !D $
\RightLabel{\LeftR{\Weakening}}
\UnaryInf$ !C, !D \fCenter !D$
\Axiom$ !C \fCenter !C $
\RightLabel{\LeftR{\Weakening}}
\UnaryInf$ !D, !C \fCenter !C$
\RightLabel{\LeftR{\Exchange}}
\UnaryInf$ !C, !D \fCenter !C$
\RightLabel{\RightR{\land}}
\BinaryInf$ !C, !D \fCenter !D \land !C $
\end{prooftree}
\end{ex}

\end{document}
